\chapter{应用验证与讨论}\label{chap:application}
\section{应用场景}
应用章节需要说明场景与前文实验之间的联系，并区分定量精度验证和可视化展示。场景图可以说明覆盖范围及几何结构，但不能替代位置误差统计。

\section{存储规模统计}
表\ref{tab:storage}展示不同单位和小数位的写法。相同指标在正文和表格中的名称、单位及舍入方式应保持一致。对已舍入数值再次舍入可能产生末位冲突，应直接使用原始统计进行显示格式转换。

\begin{table}[htbp]
  \centering
  \caption{地图存储规模格式示例}\label{tab:storage}
  \begin{tabular}{@{}lrrr@{}}
    \toprule
    序列 & 里程（km） & 点云地图（MB） & 轻量地图（MB）\\
    \midrule
    序列1 & 1.2345 & 128.25 & 0.0123\\
    序列2 & 2.3456 & 256.50 & 0.0234\\
    \bottomrule
  \end{tabular}
  \tablenote{注：仅演示数值右对齐与小数位，不代表真实文件大小。}
\end{table}

\section{本章小结}
本章给出应用讨论与存储表格的排版示例。实际论文中的指标范围、统计对象和数据来源需要与实验材料对应。
